\begin{abstract}
At Hurst index $H=0.9$, the short-memory cutoff for a nominal $5\,\%$ Bartlett self-normalised test of a
time-series mean gives a limiting rejection rate of \pthree{naive-sn-0.90-0.05}\,\% under the null.
\CosCert{} computes critical values for this test and for the fixed-$m$ periodogram and equal-weighted cosine tests.
Eigenvalue bounds for the fractional Brownian bridge and adaptive quadrature in complex ball arithmetic
give deterministic enclosures for Bartlett probabilities. We certify all \pthree{e2q-n-brackets} tabulated Bartlett
critical values across six levels and \pthree{e2q-n-H} values of $H$. The $5\,\%$ critical value rises from
\pthree{sn-q-0.50-0.950} at $H=0.5$ to \pthree{sn-q-0.99-0.950} at $H=0.99$.
At $H=0.75$, one probability evaluation near the 95th percentile takes \pthree{e5-T4-cert-s}\,s including setup,
a \pthree{e5-T4-speedup}-fold speedup over the midpoint benchmark.
Rescaling Student-$t$ cutoffs keeps the limiting size of the fixed-smoothing tests within
\pthree{rs-maxerr-090-0.05} percentage points of $5\,\%$ on the grid through $H=0.9$.
In simulations with estimated memory, plug-in tests reject up to \pthree{pl-pure-max-100}\,\% of the time at $n=100$
under pure long memory. For the four-frequency periodogram test, using an upper confidence bound for $H$ gives
rejection rates at most \pthree{comp-hi4-up-n500-max}\,\% in the $n=500$ designs.
For the Nile minima, calibrated intervals are \pthree{emp-ratio-min}--\pthree{emp-ratio-max} times as wide as
short-memory intervals, whose coverage under the fitted Gaussian model falls to \pthree{emp-cover-nv-min}\,\%.
Of \pthree{audit-n-values} published critical values, \pthree{audit-n-AB} yield rejection rates within three
Monte Carlo standard errors of nominal size.
\end{abstract}

\begin{keyword}
fractional Brownian motion \sep long memory \sep self-normalisation \sep fixed-$b$ asymptotics \sep critical values
\sep certified computation \sep generalised chi-square distribution
\MSC[2020] 62M10 \sep 60G22 \sep 62E15 \sep 65G30
\end{keyword}
